\documentclass[10pt,a4paper]{amsart}
\usepackage{amsmath,amssymb,amsthm,graphicx}
\usepackage{verbatim}
\newenvironment{reference}{\comment}{\endcomment}
%\newenvironment{reference}{}{}
\newenvironment{slogan}{\comment}{\endcomment}
\newenvironment{history}{\comment}{\endcomment}

% For commutative diagrams we use Xy-pic
\usepackage[all]{xy}

% We use 2cell for 2-commutative diagrams.
\xyoption{2cell}
% Theorem environments.
%
\theoremstyle{plain}
\newtheorem{theorem}[subsection]{Theorem}
\newtheorem{proposition}[subsection]{Proposition}
\newtheorem{lemma}[subsection]{Lemma}

\theoremstyle{definition}
\newtheorem{definition}[subsection]{Definition}
\newtheorem{example}[subsection]{Example}
\newtheorem{exercise}[subsection]{Exercise}
\newtheorem{situation}[subsection]{Situation}

\theoremstyle{remark}
\newtheorem{remark}[subsection]{Remark}
\newtheorem{remarks}[subsection]{Remarks}

\numberwithin{equation}{subsection}

% Macros
%
\def\lim{\mathop{\mathrm{lim}}\nolimits}
\def\colim{\mathop{\mathrm{colim}}\nolimits}
\def\Spec{\mathop{\mathrm{Spec}}}
\def\Hom{\mathop{\mathrm{Hom}}\nolimits}
\def\Ext{\mathop{\mathrm{Ext}}\nolimits}
\def\SheafHom{\mathop{\mathcal{H}\!\mathit{om}}\nolimits}
\def\SheafExt{\mathop{\mathcal{E}\!\mathit{xt}}\nolimits}
\def\Sch{\mathit{Sch}}
\def\Mor{\mathop{\mathrm{Mor}}\nolimits}
\def\Ob{\mathop{\mathrm{Ob}}\nolimits}
\def\Sh{\mathop{\mathit{Sh}}\nolimits}
\def\NL{\mathop{N\!L}\nolimits}
\def\CH{\mathop{\mathrm{CH}}\nolimits}
\def\proetale{{pro\text{-}\acute{e}tale}}
\def\etale{{\acute{e}tale}}
\def\QCoh{\mathit{QCoh}}
\def\Ker{\mathop{\mathrm{Ker}}}
\def\Im{\mathop{\mathrm{Im}}}
\def\Coker{\mathop{\mathrm{Coker}}}
\def\Coim{\mathop{\mathrm{Coim}}}

% Boxtimes
%
\DeclareMathSymbol{\boxtimes}{\mathbin}{AMSa}{"02}

%
% Macros for moduli stacks/spaces
%
\def\QCohstack{\mathcal{QC}\!\mathit{oh}}
\def\Cohstack{\mathcal{C}\!\mathit{oh}}
\def\Spacesstack{\mathcal{S}\!\mathit{paces}}
\def\Quotfunctor{\mathrm{Quot}}
\def\Hilbfunctor{\mathrm{Hilb}}
\def\Curvesstack{\mathcal{C}\!\mathit{urves}}
\def\Polarizedstack{\mathcal{P}\!\mathit{olarized}}
\def\Complexesstack{\mathcal{C}\!\mathit{omplexes}}
% \Pic is the operator that assigns to X its picard group, usage \Pic(X)
% \Picardstack_{X/B} denotes the Picard stack of X over B
% \Picardfunctor_{X/B} denotes the Picard functor of X over B
\def\Pic{\mathop{\mathrm{Pic}}\nolimits}
\def\Picardstack{\mathcal{P}\!\mathit{ic}}
\def\Picardfunctor{\mathrm{Pic}}
\def\Deformationcategory{\mathcal{D}\!\mathit{ef}}
\usepackage{cleveref}
\crefname{originaldefinition}{original definition}{original definitions}
\Crefname{originaldefinition}{Original definition}{Original definitions}
\makeatletter
\newlabel{definition-quotient-stack}{{044Q}{source}{}{sourcepointer-1}{}}
\newlabel{definition-quotient-stack@cref}{{[originaldefinition][0][]044Q}{[1][0][]source}{}{}{}}
\makeatother
\crefname{originaldefinition}{original definition}{original definitions}
\Crefname{originaldefinition}{Original definition}{Original definitions}
\makeatletter
\newlabel{definition-stabilizer-groupoid}{{0448}{source}{}{sourcepointer-2}{}}
\newlabel{definition-stabilizer-groupoid@cref}{{[originaldefinition][0][]0448}{[1][0][]source}{}{}{}}
\makeatother
\crefname{originalreference}{original reference}{original references}
\Crefname{originalreference}{Original reference}{Original references}
\makeatletter
\newlabel{stacks-lemma-stackification-inertia}{{06NS}{source}{}{sourcepointer-3}{}}
\newlabel{stacks-lemma-stackification-inertia@cref}{{[originalreference][0][]06NS}{[1][0][]source}{}{}{}}
\makeatother
\crefname{originalreference}{original reference}{original references}
\Crefname{originalreference}{Original reference}{Original references}
\makeatletter
\newlabel{categories-lemma-inertia-fibred-category}{{034H}{source}{}{sourcepointer-4}{}}
\newlabel{categories-lemma-inertia-fibred-category@cref}{{[originalreference][0][]034H}{[1][0][]source}{}{}{}}
\makeatother
\title{The inertia stack of a quotient stack is a quotient of the stabilizer}
\author{The Stacks Project Authors}
\date{stacks-project:spaces-groupoids.tex:local-snapshot-20261012; GNU FDL 1.2 or later}
\begin{document}
\maketitle
\noindent\textit{Editorial note.} All delivered mathematical text is the author's own, copied verbatim from the pinned Stacks Project chapter file spaces-groupoids.tex; no statement or proof step was rewritten, reordered or removed.\par\medskip
\noindent\textit{Editorial note.} Span recomputation by label pairing (LANE\_BRIEF 10.2): in this chapter the author's proof environment follows the statement environment as a standalone block; the spec covers the complete environment pair (statement span, proof span) as recorded in the pool candidate.\par\medskip
\ifdefined\hypertarget\hypertarget{sourcepointer-1}{}\fi
\noindent\textit{Source reference.} Original definition 044Q: sources/stacks-chow/tags.txt, source lines 5409--5409. Cross-chapter standard prerequisite cited by the author inside the counted span; it is used as a pointer and its own proof is not part of this item.\par\medskip
\ifdefined\hypertarget\hypertarget{sourcepointer-2}{}\fi
\noindent\textit{Source reference.} Original definition 0448: sources/stacks-chow/tags.txt, source lines 5391--5391. Cross-chapter standard prerequisite cited by the author inside the counted span; it is used as a pointer and its own proof is not part of this item.\par\medskip
\ifdefined\hypertarget\hypertarget{sourcepointer-3}{}\fi
\noindent\textit{Source reference.} Original reference 06NS: sources/stacks-chow/tags.txt, source lines 8546--8546. Cross-chapter standard prerequisite cited by the author inside the counted span; it is used as a pointer and its own proof is not part of this item.\par\medskip
\ifdefined\hypertarget\hypertarget{sourcepointer-4}{}\fi
\noindent\textit{Source reference.} Original reference 034H: sources/stacks-chow/tags.txt, source lines 4092--4092. Cross-chapter standard prerequisite cited by the author inside the counted span; it is used as a pointer and its own proof is not part of this item.\par\medskip
\begin{lemma}
\label{lemma-presentation-inertia}
Assume $B \to S$ and $(U, R, s, t, c)$ as in
Definition \ref{definition-quotient-stack} (1).
Let $G/U$ be the stabilizer group algebraic space of the groupoid
$(U, R, s, t, c, e, i)$, see
Definition \ref{definition-stabilizer-groupoid}.
Set $R' = R \times_{s, U} G$ and set
\begin{enumerate}
\item $s' : R' \to G$, $(r, g) \mapsto g$,
\item $t' : R' \to G$, $(r, g) \mapsto c(r, c(g, i(r)))$,
\item $c' : R' \times_{s', G, t'} R' \to R'$,
$((r_1, g_1), (r_2, g_2) \mapsto (c(r_1, r_2), g_1)$.
\end{enumerate}
Then $(G, R', s', t', c')$ is a groupoid in algebraic spaces over $B$
and
$$
\mathcal{I}_{[U/R]} = [G/ R'].
$$
i.e., the associated quotient stack is the inertia stack of $[U/R]$.
\end{lemma}

\begin{proof}
By
Stacks, Lemma \ref{stacks-lemma-stackification-inertia}
it suffices to prove that $\mathcal{I}_{[U/_{\!p}R]} = [G/_{\!p} R']$.
Let $T$ be a scheme over $S$. Recall that an object of the inertia fibred
category of $[U/_{\!p}R]$ over $T$ is given by a pair
$(x, g)$ where $x$ is an object of $[U/_{\!\!p}R]$ over $T$
and $g$ is an automorphism of $x$ in its fibre category over $T$.
In other words, $x : T \to U$ and $g : T \to R$ such that
$x = s \circ g = t \circ g$. This means exactly that
$g : T \to G$. A morphism in the inertia fibred category
from $(x, g) \to (y, h)$ over $T$ is given by
$r : T \to R$ such that $s(r) = x$, $t(r) = y$
and $c(r, g) = c(h, r)$, see the commutative diagram in
Categories, Lemma \ref{categories-lemma-inertia-fibred-category}.
In a formula
$$
h = c(r, c(g, i(r))) = c(c(r, g), i(r)).
$$
The notation $s(r)$, etc is a short hand for $s \circ r$, etc.
The composition of $r_1 : (x_2, g_2) \to (x_1, g_1)$
and $r_2 : (x_1, g_1) \to (x_2, g_2)$ is
$c(r_1, r_2) : (x_1, g_1) \to (x_3, g_3)$.

\medskip\noindent
Note that in the above we could have written $g$ in stead of $(x, g)$
for an object of $\mathcal{I}_{[U/_{\!p}R]}$ over $T$ as $x$ is the
image of $g$ under the structure morphism $G \to U$. Then the morphisms
$g \to h$ in $\mathcal{I}_{[U/_{\!p}R]}$ over $T$ correspond exactly
to morphisms $r' : T \to R'$ with $s'(r') = g$ and $t'(r') = h$.
Moreover, the composition corresponds to the rule explained in (3).
Thus the lemma is proved.
\end{proof}

\end{document}
